\section{Conclusion}

We have demonstrated that the conserved branching exponent $\alpha^*$ observed
in biological transport networks is not merely an empirical curiosity, but a
\chA{structural consequence} of scale-free evolution. The transition from local
junctional control to global network coordination is \chA{driven by} the fundamental
incommensurability of the physical costs governing the vascular tree.

The logical edifice of the incommensurability principle rests on three
interlocking pillars. First, the \textbf{informational cost of local
optimization} (Proposition~\ref{thm:nogo}) reveals that for a cell to maintain
architectural universality through local sensing, it would need to measure and
compensate for non-local invariants with a precision that exceeds the capacity
of molecular receptors. The vascular tree avoids this cost by converging to the
\emph{network-level minimax}.

Second, the functional form of this attractor is uniquely dictated by
\textbf{Metabolic Scaling Symmetry}. By grounding the fitness landscape in the
ATP stoichiometry of metabolism, we have shown that the linear fractional excess
is the \chA{admissible cost functional under the stated thermodynamic and compositional assumptions}.

Third, the transition between physical regimes is governed by the
\textbf{Kinematic Matching Criterion} (Theorem~\ref{thm:kinematic_matching}).
The topologically grounded dual-threshold framework---with fluid threshold
$\mathrm{Wo}_c^{\mathrm{fluid}} = \sqrt{3}$ and wave threshold
$\mathrm{Wo}_c^{\mathrm{wave}} = 3/\sqrt{2}$---is derived from the evanescent
nature of the transverse velocity component at the bifurcation junction,
combined with the Navier-Stokes expression for viscous absorption capacity
($\mathcal{Q}^{-1} = 6/\mathrm{Wo}^2$). The critical condition $\mathcal{Q}^{-1}
= d-1$---the vascular analogue of the Ioffe-Regel criterion---marks the
threshold where geometric scattering overwhelms viscous damping, establishing
the onset of coherent wave propagation. The agreement with the observed
allometric transition in mammals \chA{supports} the claim that the fundamental asymptotics of
fluid kinematics constrain the coarse-grained architecture of cardiovascular
systems.

The incommensurability principle thus provides a unified framework that subsumes
Murray's Law and the area-preserving limit as the two viscous and inertial
boundaries of a single minimax landscape. Beyond vascular biology, this
principle applies to any adaptive network forced to reconcile physically
distinct cost dimensions.

\chA{As shown by} the Topological Rigidity corollary
(Corollary~\ref{thm:rigidity}), the minimax saddle point requires no target-exponent
fine-tuning: the branching exponent $\alpha^*$ depends only on the dimensionless
structural parameters $(G, N, p, \alpha_w, A)$ and is completely independent of
all pure metabolic parameters---blood oxygen cost, proximal flow, segment
length, ATP stoichiometry, cardiac output ($|S_x| < 0.01$,
Table~\ref{tab:sensitivity})---while exhibiting only weak residual sensitivity
to fluid-mechanical parameters (viscosity, wall metabolism: $|S_x| < 0.2$). The
architecture is effectively decoupled from the biochemistry.

\begin{table}[H]
\centering
\caption{Separation of genuine predictions from fitted parameters}
\begin{tabular}{lccc}
\toprule
Quantity & Value & Epistemological Status & Source \\
\midrule
\multicolumn{4}{l}{\textbf{\chA{Derived outputs (no target-exponent fit)}}} \\
$\mathrm{Wo}_c$ (3D fluid) & $\sqrt{3} \approx 1.732$ & \chA{Physical criterion; exact-Bessel checked} & Eq.~\eqref{eq:woc_closed_form} \\
$\mathrm{Wo}_c$ (2D fluid) & $\sqrt{6} \approx 2.449$ & \chA{Physical criterion; exact-Bessel checked} & Eq.~\eqref{eq:woc_closed_form} \\
 $M^*$ (transition mass) & $0.8$\,g & \chA{Model-derived transition landmark} & Eq.~\eqref{eq:wo_scaling} \\
$\alpha^*$ (symmetric) & $\chA{\VarAlphaStarModel}$ & Minimax solution & Eq.~\eqref{eq:eta_star} \\
\midrule
\multicolumn{4}{l}{\textbf{Empirical Inputs (measured independently)}} \\
$p_{eff}$ (wall thickness) & $0.77$ & Derived effective exponent & Guo 2003 \cite{Guo2003} \\
$\alpha^*$ (empirical target) & $2.72$ & Morphometric measurement & Kassab et al. 1993 \\
\midrule
\multicolumn{4}{l}{\textbf{Postdictions (explained by measured heterogeneity)}} \\
$\Delta\alpha$ (residual) & $0.094$ & Asymmetry + taper & Appendix C \\
\bottomrule
\end{tabular}
\end{table}

\subsection{A Critical Falsifiable Test: Ontogenetic Phase Transition}

The theory makes a sharp, falsifiable prediction regarding vascular
morphogenesis. In early embryonic development, the aortic radius is sufficiently
small that the Womersley number satisfies $\mathrm{Wo} \ll
\mathrm{Wo}_c^{\mathrm{fluid}}$ throughout the nascent arterial tree. Under
these conditions, wave reflections are negligible, and the network should
spontaneously adopt Murray's viscous-optimal scaling $\alpha \approx 3.0$. As
the organism grows and the aortic Womersley number crosses the critical fluid
threshold $\mathrm{Wo}_0(M^*) = \sqrt{3}$ (corresponding to a body mass $M^*
\approx \VarMStar\,$g in mammals), the theory predicts an abrupt structural
remodeling: the branching exponent must shift from the viscous attractor
($\alpha \approx 3.0$) to the wave-influenced minimax attractor ($\alpha^*
\approx \VarAlphaStar$).

This prediction is directly testable via time-resolved morphometry of developing
embryos. Two-photon microscopy of zebrafish embryos (transparent, rapid
development) or synchrotron micro-CT of staged mouse embryos can measure the
generation-by-generation branching exponent $\alpha(t)$ as a function of
developmental time (and thus body mass). The critical measurement is whether
$\alpha$ exhibits a systematic transition from $\sim 3.0$ at early stages (when
$M < M^*$) to $\sim \VarAlphaStar$ at later stages (when $M > M^*$).

\textbf{Falsification Criterion:} If embryonic vasculature exhibits $\alpha
\approx \VarAlphaStar$ from the earliest stages of angiogenesis---when wave
reflections are physically absent---the theory is falsified. Such an observation
would demonstrate that the branching exponent is a genetically hard-coded
blueprint rather than an emergent hydrodynamic attractor, invalidating both the
Epistemic Bound (Proposition~\ref{thm:nogo}) and the proposed need for network-level
minimax optimization. Conversely, observation of the predicted viscous-to-wave
transition would constitute direct experimental \chA{support} for the claim that vascular
geometry is sculpted by the incommensurability principle.

\begin{table}[H]
\centering
\caption{\textbf{Falsifiable Prediction: Ontogenetic Transition of $\alpha$}
\emph{No experimental data currently available; this constitutes a direct test
of the incommensurability principle.}}
\label{tab:ontogenetic}
\begin{tabular}{lcccl}
\toprule
Stage & Mass $M$ (g) & $\mathrm{Wo}_0$ & $\alpha_{\mathrm{pred}}$ & Regime \\
\midrule
E10 (early embryo) & 0.1  & 0.5  & \chA{2.90} & Viscous \\
E15                & 0.5  & 1.0  & 2.85 & Transition onset \\
E20                & 2.0  & 1.8  & 2.70 & Minimax attractor \\
Neonate            & 25   & 3.5  & 2.70 & Minimax (stable) \\
\bottomrule
\end{tabular}
\end{table}

\chA{The masses and Womersley numbers above trace a single organism's
\emph{ontogenetic} (intraspecific) growth trajectory. Along this trajectory
heart rate, cardiac output, and aortic calibre scale with developmental mass
differently from the \emph{interspecific} adult allometry
($\dot{Q}\propto M^{3/4}$, $f\propto M^{-1/4}$) that yields
$\mathrm{Wo}_0\propto M^{1/4}$ and the cross-species transition mass
$M^*\approx\VarMStar\,$g (\S\ref{sec:womersley_minimax}). The tabulated
developmental schedule therefore rises more steeply than $M^{1/4}$
(here $\mathrm{Wo}_0$ grows roughly as $M^{0.4}$ over the embryonic range), and
the stage at which a given species crosses $\mathrm{Wo}_0=\sqrt{3}$
($\approx 2\,$g in this illustrative schedule) is trajectory-specific and is
\emph{not} expected to coincide with the interspecific $M^*$. The robust,
falsifiable content is qualitative: the branching exponent must migrate from the
viscous value ($\alpha\approx 3.0$) toward the minimax attractor
($\alpha^*\approx\VarAlphaStar$) as $\mathrm{Wo}_0$ crosses its critical
threshold, whatever the precise developmental mass at which this occurs.}

The convergence of evolution toward a scale-free attractor that emerges from the
irreconcilability of physical costs is not unique to biology. Wherever adaptive
systems must navigate incommensurable constraints---whether in neural networks
optimizing speed versus accuracy, or ecological food webs balancing energy
acquisition versus predation risk---the minimax principle offers a universal
resolution. Universality, in this view, does not require fine-tuning; it
requires only that the costs be incommensurable.
